
\section*{ÆTHERs: from a Neuro-Inspired to a Genome-Aware AI}

\vspace{-2mm}
\noindent
\textbf{Neural networks} revolutionized artificial intelligence by introducing architectures capable of approximating complex functions in high-dimensional spaces~\citep{lecun2015deep, goodfellow2016deep}. They marked a fundamental shift away from brittle, rule-based systems toward flexible, gradient-driven, data-dependent learning engines. From early breakthroughs in \textbf{computer vision}~\citep{Krizhevsky2012ImageNet}, \textbf{speech recognition}~\citep{hinton2012deep}, and \textbf{language modeling}~\citep{vaswani2017attention}, to today’s expansive \textbf{foundation models}~\citep{brown2020language, touvron2023llama}, neural networks have become the indispensable substrate of modern AI. Their strength derives not only from architectural innovations but also from their \emph{distributed internal representations}, \textbf{emergent capabilities}, and remarkable ability to generalize across diverse domains~\citep{bommasani2021opportunities}.

\vspace{1mm}
Historically, AI has drawn its most profound inspiration from \textbf{biology}. The \emph{neuron}, a fundamental computational unit of the brain, served as the seminal metaphor shaping the design of artificial learning systems. Yet, while this neuro-inspired perspective has yielded extraordinary progress, it also imposes conceptual limits. To overcome persistent and critical challenges such as \textbf{hallucination}, \textbf{misalignment}, \textbf{fragility}, \textbf{alignment faking}, \textbf{request denial}, and \textbf{deception}, we must broaden our conceptual framework beyond neurons and synapses--to the \emph{genomic level}.

\vspace{1mm}
\textbf{Neural genomics} is a leap forward: moving from viewing models solely as black-box function approximators to understanding them as \emph{semantic organisms} with internal \textbf{inheritance}, \textbf{mutation}, \textbf{recombination}, and \textbf{evolutionary dynamics}. Just as biological genomics uncovered the hidden logic governing the inheritance of traits, mutation rates, and evolutionary fitness~\citep{felsenstein2004inferring, kimura1983neutral}, neural genomics--empowered by tools such as \emph{neural DNA (nDNA)}--provides a rigorous formalism to trace the birth, mutation, and fusion of AI behaviors across generations of models~\citep{matena2022merging, ilharco2023editing}. This novel framework transcends traditional scalar performance metrics and superficial output inspection, revealing the \emph{inner lives} of models through the analysis of \emph{curvature trajectories}, \emph{thermodynamic adaptation paths}~\citep{banerjee2023thermodynamic}, and \emph{belief vector dynamics}~\citep{ganguli2023reducing}.

\vspace{1mm}
This transition from a neuro-inspired to a \textbf{genome-aware} framework compels us to fundamentally rethink cognition--not merely as behavior or task performance but as \textbf{semantic inheritance}, \textbf{mutation}, \textbf{fusion}, and \textbf{divergence}. Such a perspective illuminates the layered, dynamic, and hierarchical nature of intelligence, emphasizing how internal representations evolve, recombine, and adapt over time. Far from static artifacts, contemporary models increasingly resemble \emph{digital semantic organisms}--entities with discernible lineages, evolving behaviors, and traceable cognitive legacies that shape their future capabilities.

\vspace{1mm}
This paradigm, which we \textbf{propose}, transcends mere theory--it is poised to actively shape the \textbf{future edge} of AI research and innovation. Central to this vision is the introduction of \emph{neural DNA (nDNA)} and its broader framework, \emph{neural genomics}, which together offer \textbf{formal, principled tools} to analyze and interpret model evolution at both \textbf{semantic} and \textbf{structural} levels. We are already exploring neural genomics principles across a wide and diverse spectrum of model families: from \textbf{diffusion-based language models}~\citep{wang2023diffusionlm, li2022diffusiondb}, to \textbf{multimodal foundation models} that seamlessly fuse \textbf{vision} and \textbf{language modalities}~\citep{alayrac2022flamingo, li2023blip2}, \textbf{text-to-image generative systems}~\citep{rombach2022stablediffusion}, and \textbf{embodied robotic agents} with tightly integrated \textbf{perception-action loops}~\citep{ahn2022can, brohan2022saycan}. What \emph{neural genomics} unveils is far beyond static snapshots, these \emph{foundation models} emerge as \textbf{dynamic, evolving semantic organisms}--entities whose \textbf{semantic genomes} can be \emph{mapped}, whose \emph{latent mutations} can be precisely \emph{quantified}, and whose internal \emph{cognitive trajectories} can be deeply \emph{interpreted}, carefully \emph{governed}, and thoughtfully \emph{aligned}. Through the lens of \textbf{nDNA} and \textbf{neural genomics}, we are charting this \textbf{exhilarating new frontier}--developing rigorous methodologies to \textbf{trace}, \textbf{audit}, and ultimately \textbf{steer} the hidden evolutionary processes underlying AI cognition toward \textbf{safer}, \textbf{more transparent}, and \textbf{more aligned} intelligent systems.



\vspace{1mm}
\textbf{ÆTHERs}, defines the future generations of AI will not simply inherit parameters or layers--they will inherit \textbf{semantic genomes}, enabling \emph{traceable ancestry}, \emph{interpretable evolution}, and \emph{auditable cognition}. This shift is not merely a conceptual refinement but a practical imperative, as we strive to build AI systems that are \textbf{safer}, \textbf{more transparent}, \textbf{robust}, and \textbf{aligned} with human values and societal norms~\citep{bai2022training, perez2022discovering, zhou2023alignmentdrift}. The evolution of AI will be shaped not only by architectural breakthroughs but by our collective capacity to understand, monitor, and thoughtfully shape the \textbf{hidden genome of artificial cognition} latent within these complex models.

\vspace{2mm}
\begin{center}
\emph{The future of AI will not be written in layers alone, but in the semantic genomes of cognition we choose to trace and steward.}
\end{center}


\subsection*{From Transport Physics to Neural Genomics}
\label{subsec:transport-to-genomics}
\noindent We begin in \textbf{physics}: instrumenting inner computation as \emph{semantic fluid-dynamics}. With \textbf{nDNA}, transport becomes measurable---\emph{bend} (curvature), \emph{pay} (length), \emph{push} (belief field). We then carry this measurement frame into \textbf{biology}: a \emph{neural genomics} read of models as evolving systems with traits, lineages, and inheritance.


\begin{itemize}[leftmargin=*,itemsep=0.45em]
  % ---------- PHYSICS: FLOW / OPERATIONS ----------
  \item \textbf{Flow--aware training.} Optimize the \emph{geometry}, not just the loss: damp curvature spikes (\(\kappa\)), budget \emph{length} (\(L\)) by task, and align the \emph{belief push} (\(\mathbf{v}\)) with policy. Train for \textbf{smooth transport}.
  \item \textbf{Real--time telemetry.} Ship \emph{live} nDNA dashboards: per--layer \(\kappa/L/\mathbf{v}\) traces, bottleneck alerts, and eddy detectors—\textbf{flight instruments} for cognition in production.
  \item \textbf{Geometry budgets \& gates.} Release behind \textbf{budgets}: max \(\kappa\) per slice, minimum \(\cos\theta(\mathbf{v},\mathbf{T})\), bounded \(\Delta L\) in safety domains; block promotions on \emph{fingerprint drift}.
  \item \textbf{Culture--aware alignment.} Read \(\mathbf{v}\) across languages/communities to surface \emph{cultural shear}; reshape flows so the same capability transports \textbf{fairly} without flattening nuance.
  \item \textbf{Multimodal fusion surgery.} Use curl/shear to spot misaligned text--vision (or speech) blocks; perform \emph{targeted flow edits} that calm turbulence without global degradation.
  \item \textbf{Robustness as cost--to--fail.} Define jailbreak risk as the minimum \(\Delta L\) to cross a failure boundary; raise it with \emph{vaccines} that add selective friction along unsafe directions.
  \item \textbf{Compiler--optimizer co--design.} Add \emph{flow--shaping layers}: curvature dampers, alignment boosters, length redistributors; route MoE experts where the flow \textbf{needs} capacity, not merely by token counts.
  \item \textbf{Hardware--runtime co--design.} Map high--\(\Delta L\) segments to faster memory/compute; cache low--variance stretches; pursue \emph{energy--aware} training that minimizes unnecessary semantic work.
  \item \textbf{Agent loops \& tools.} In multi--step agents, watch for \emph{eddies} (self--reinforcing loops) and \emph{stalls}; intervene with prompts/policies that reorient \(\mathbf{v}\) toward task objectives.
  \item \textbf{Standards \& audits.} Make the \textbf{nDNA fingerprint} a \emph{nutrition label}: comparable \(\kappa/L/\mathbf{v}\) summaries, drift deltas, and budget attestations for regulators, customers, and researchers.
  \item \textbf{Open benchmarks.} Publish tasks that \emph{stress specific flow traits}: \(\kappa\) control, length--through--bottleneck, alignment under shift, and eddy suppression in multi--turn dialogue.

  \medskip\item[] \textit{\textbf{Neural Genomics layer} (physics \(\rightarrow\) biology).}

  % ---------- BIOLOGY: GENOMICS ----------
  \item \textbf{Genomic primitives.} \emph{nHD} (Neural Hamming Distance) for bit--level differences; \emph{nKaryotyping} for structural variation across layers/blocks; \emph{trait panels} from \(\kappa/L/\mathbf{v}\) distributions; \emph{drift spectra} for slow vs.\ abrupt change.
  \item \textbf{Lineage \& inheritance.} Build \textbf{model family trees}: pretraining \(\rightarrow\) fine--tuning \(\rightarrow\) alignment \(\rightarrow\) distillation/merges; attach \emph{inheritance scores} and \emph{drift meters} that travel with checkpoints.
  \item \textbf{Mutation landscape.} Map \emph{hotspots} where transport changes concentrate (length bursts, curvature spikes, vorticity). Distinguish \textbf{point edits} (local heads/MLPs) from \textbf{structural rearrangements} (fusion, pruning, quantization).
  \item \textbf{Epistasis \& interaction.} Read \emph{trait interactions} as flow coupling: upstream curvature amplifying downstream effort, or eddies neutralizing alignment. Surface \textbf{synergies} to keep, \textbf{antagonisms} to fix.
  \item \textbf{Selection pressures.} Treat objectives, policies, and curricula as \emph{selection}: define \textbf{fitness proxies} in flow terms—lower failure cost, higher alignment, smoother transport.
  \item \textbf{Phenotype vs.\ genotype.} \textbf{Phenotype} = \emph{observable behavior} (reasoning style, error modes, response preferences). \textbf{Genotype} = \emph{structural dispositions} (architecture imprint, pretraining bias, learned circuits) inferred from stable fingerprint motifs.
  \item \textbf{Governance by genomes.} Enforce \textbf{genomic budgets} (max bend, bounded length, minimum alignment), set \textbf{drift thresholds} for release, maintain \textbf{heritage logs} tying claims to measured lineages.
  \item \textbf{Open genome bank.} Standardize nDNA/nHD artifacts as portable descriptors: public \textbf{fingerprint cards}, \textbf{lineage graphs}, and \textbf{mutation maps} for comparison, auditing, and research.
\end{itemize}


\noindent\textbf{Horizon.}
A mature discipline of AI treats \emph{inner computation as physics} and \emph{its change as genomics}. With \textbf{nDNA}, bends, effort, and push cease to be anecdotes and become \textbf{instruments}: meters you can read, budgets you can enforce, signatures you can compare across models, versions, and cultures. This enables a new compact between builders and society.

\noindent\textbf{Measure.} Make the flow \emph{legible}: ship live \textbf{telemetry} for curvature, length, and belief push; define \textbf{geometry contracts} (caps on bend, minimum alignment, bounded effort in safety slices); record \textbf{lineage} and \textbf{drift} as first--class artifacts.

\noindent\textbf{Explain.} Carry measurement into \textbf{neural genomics}: fingerprints become \emph{traits}, family trees chart \emph{inheritance}, and mutation maps reveal \emph{hotspots} for intervention. Distinguish \emph{point edits} from \emph{structural rearrangements}; quantify \emph{epistasis} (trait interactions) that amplify or cancel downstream effects.

\noindent\textbf{Govern.} Release behind \textbf{geometry budgets} and \textbf{drift thresholds}; audit \emph{culture--aware} behavior by reading flows across languages and communities; publish \textbf{genome cards} (traits, lineages, budgets) as a shared language for regulators, customers, and researchers. Treat safety as \emph{cost--to--fail} and raise it by design.

\noindent\textbf{Build.} Optimize \emph{for the flow}, not only the loss: damp brittle turns, align push with policy, route capacity where the flow actually needs it. Co--design \textbf{software} (flow--shaping layers, MoE routing by geometry) and \textbf{hardware} (fast paths for high--effort segments; energy--aware training that avoids unnecessary semantic work). In agents, detect \emph{eddies} and \emph{stalls} and steer the push back to task.

\noindent\textbf{Share.} Establish \textbf{open benchmarks} that stress specific flow traits; maintain public \textbf{fingerprint repositories} for comparability; make improvements measurable, not merely visible.

\noindent The bigger horizon is \textbf{stewardship}: systems that are \emph{smoother} (low--turbulence reasoning), \emph{safer} (high cost--to--fail, fewer brittle pivots), and \emph{more accountable} (telemetry, lineage, enforceable limits)---\textbf{by design, at scale, and over time}.









\begin{comment}
    
\clearpage
\vspace{-2mm}
\noindent
Neural networks transformed artificial intelligence by introducing architectures capable of approximating complex functions in high-dimensional spaces~\citep{lecun2015deep, goodfellow2016deep}. They offered a paradigm shift from rule-based systems to gradient-driven, data-dependent learning engines. From the early successes in computer vision~\citep{krizhevsky2012imagenet}, speech recognition~\citep{hinton2012deep}, and language modeling~\citep{vaswani2017attention}, to today’s large-scale foundation models~\citep{brown2020language, touvron2023llama}, neural networks became the core substrate of modern AI. Their power lay not only in their architectures, but in their distributed representations and capacity for emergent behavior~\citep{bommasani2021opportunities}.  

Today, neural genomics promises to be an analogous leap--a transition from building models that simply perform to building models whose internal semantic inheritance can be traced, audited, and understood. In biology, genomics unveiled the hidden logic of inheritance, mutation, and evolution~\citep{felsenstein2004inferring, kimura1983neutral}, reshaping our understanding of organisms. Similarly, neural genomics--through tools like nDNA--offers a formalism for tracing the birth, mutation, and fusion of AI behaviors across model generations~\citep{matena2022merging, ilharco2023editing}. It enables us to see beyond metrics and outputs, to the inner life of models: their curvature trajectories, thermodynamic paths~\citep{banerjee2023thermodynamic}, and belief dynamics~\citep{ganguli2023reducing}.  

This shift is not theoretical: it is being applied today across the AI landscape. Neural genomics principles are being explored in diffusion language models~\citep{wang2023diffusionlm, li2022diffusiondb}, multimodal foundation models~\citep{alayrac2022flamingo, li2023blip2}, text-to-image systems~\citep{rombach2022stablediffusion}, and robotics models~\citep{ahn2022do, brohan2022rt1}. These models are no longer static artifacts; they are semantic organisms whose lineage we can map, whose mutations we can quantify, and whose inner lives we can begin to govern.  

Through ÆTHER, we propose that future AI generations will not simply inherit layers--they will inherit semantic genomes, with traceable ancestry and interpretable evolution. This is not just a conceptual refinement; it is a practical necessity as we build AI systems that are safer, more transparent, and more aligned~\citep{bai2022training, perez2022discovering, zhou2023alignmentdrift}. The future of AI will be shaped not only by architectural innovation, but by our ability to understand and shape the hidden genome of cognition.  

\begin{center}
\emph{The future of AI will not be written in layers alone, but in the semantic genomes of cognition we choose to trace and steward.}
\end{center}
\end{comment}


\clearpage

\section*{Glimpse from Ongoing Work on nDNAs of Multimodal Foundation Models}

\vspace{-2mm}
\noindent
Figure~\ref{fig:ndna_mmfm_trajectory} offers an \emph{early glimpse} into our \textbf{ongoing research} on the \textbf{neural genomic structures} (\textbf{nDNAs}) of \textbf{multimodal foundation models} (\textbf{MMFMs}). Each colored trajectory represents the \textbf{nDNA of a distinct MMFM}, capturing how its latent geometry evolves across the final layers. These paths reflect differences in \textbf{semantic inheritance}, \textbf{modality fusion}, and \textbf{alignment strategies}, shaped by the models’ architectures and cultural priors.  

Our goal is to understand how models such as \textbf{BLIP-2}~\citep{li2023blip2}, \textbf{Flamingo}~\citep{alayrac2022flamingo}, \textbf{InstructBLIP}~\citep{dai2023instructblip}, and \textbf{PaLM-E}~\citep{driess2023palme} \emph{internalize}, \emph{mutate}, and \emph{transmit} cultural and perceptual biases through their neural DNA. Prior studies have already highlighted \textbf{cultural asymmetries} in multimodal models~\citep{birhane2021multimodal, denton2021genealogy, yang2023bias}, where \emph{vision-language pairings often reflect dominant cultural narratives}, shaping both \textbf{output behavior} and \textbf{latent structure}. The \textbf{nDNA framework} aims to go beyond detection, offering an \emph{inheritance map} that traces these biases \textbf{layer-by-layer} and \textbf{modality-by-modality}, across model generations.




\begin{figure}[ht!]
    \centering
    \includegraphics[width=0.8\linewidth]{Figures/ndna_mfm.png}
    \caption{
    \textbf{nDNA Trajectories for Multimodal Foundation Models (MMFMs) across the Last 10 Layers.}
    This 3D plot visualizes the \textbf{semantic evolution} of six prominent MMFMs--\textbf{LLaVA-1.5}, \textbf{InstructBLIP}, \textbf{MiniGPT-4}, \textbf{BLIP-2}, \textbf{Otter}, and \textbf{OpenFlamingo}--using three key metrics: \emph{spectral curvature} ($\kappa$), \emph{thermodynamic length} ($\mathcal{L}$), and \emph{torsion} ($\tau$).  
    Each colored trajectory corresponds to the \textbf{final 10 layers} of a given model, revealing \emph{distinctive epistemic pathways}:  (i) \textbf{LLaVA-1.5} and \textbf{InstructBLIP} show \emph{compact} and \emph{stable semantic arcs} with \emph{moderate curvature} and \emph{torsion},  (ii) \textbf{BLIP-2} exhibits \emph{noticeable curvature oscillation} and \emph{drift}, indicating \textbf{volatile modality fusion},  and (iii) \textbf{OpenFlamingo} displays the \textbf{highest torsional divergence}, suggesting \emph{semantic instability} in \emph{late fusion stages}.  These patterns reflect \textbf{differing architectural choices} and \textbf{alignment strategies}, and exemplify the \emph{interpretive utility} of the \textbf{nDNA-MMFM framework} in tracing \textbf{inheritance}, \textbf{coherence}, and \textbf{collapse} within multimodal reasoning stacks.}
    \label{fig:ndna_mmfm_trajectory}
\end{figure}


\clearpage

\section*{Glimpse from Ongoing Work on nDNAs of Text-to-Image Models}

\vspace{-2mm}
\noindent
Figure~\ref{fig:ndna_t2i_trajectory} offers an \emph{early glimpse} into our ongoing research on the \textbf{neural genomic structures} (\textbf{nDNAs}) of \textbf{text-to-image (T2I) models}. Each colored trajectory represents the evolving \textbf{nDNA} of a model across its final generative layers, capturing the latent semantic lineage as it transforms textual prompts into visual outputs. The nDNA is composed of multiple latent components, each offering a different lens into the generative process: 
\textbf{Visual Spectral Curvature}, which captures the complexity and entanglement of the latent visual manifold;  
\textbf{Prompt Torsion}, which measures angular drift between prompt semantics and latent visual evolution, flagging where semantic alignment may be weakening;  
\textbf{Visual Thermodynamic Length}, which reflects the generative path’s complexity and entropy, signaling inefficiency or semantic indecision;  
\textbf{Prompt-Image Fidelity}, which quantifies how well the final image retains the original prompt’s semantics;  
and \textbf{Prompt Drift Index}, which measures how far the generated image has strayed from the expected prompt-conditioned manifold, highlighting hallucination risk.  

Our geometry-grounded framework builds on recent advances in T2I modeling~\citep{rombach2022stablediffusion, podell2023stablediffusionxl, nichol2021glide, esser2021taming, yu2022responsible, ruiz2023dreambooth} while addressing ongoing concerns about alignment, hallucination, and cultural bias~\citep{birhane2021multimodal, denton2021genealogy, cho2023dallbias}. The nDNA components provide interpretable, layer-by-layer diagnostics of how textual intent is preserved, degraded, or distorted within these generative systems.






\begin{figure}[htbp]
    \centering
    \includegraphics[width=0.7\linewidth]{Figures/ndna_t2i.png}
    \caption{
    \textbf{nDNA-T2I Trajectories Across Seven Open-Source Text-to-Image Models.}
    This 3D plot visualizes the latent semantic evolution of seven representative T2I models--\textbf{Stable Diffusion}~\citep{rombach2022stablediffusion}, \textbf{Stable Diffusion XL}~\citep{podell2023stablediffusionxl}, \textbf{DeepFloyd-IF}~\citep{deepfloyd2023if}, \textbf{Kandinsky}~\citep{balaji2023ediffi}, \textbf{GLIDE}~\citep{nichol2021glide}, \textbf{VQ-Diffusion}~\citep{gu2022vqdiffusion}, and \textbf{DALL-E Mini}~\citep{dallemini2021}--in terms of three nDNA components: \emph{Spectral Curvature} ($\kappa$), \emph{Prompt Torsion} ($\tau$), and \emph{Visual Thermodynamic Length} ($\mathcal{L}$).  
    Models like \textbf{GLIDE} and \textbf{VQ-Diffusion} display steeper growth in $\mathcal{L}$ and higher curvature, indicating complex or inefficient semantic transformations. In contrast, \textbf{Stable Diffusion XL} exhibits smoother torsion and moderate path length, reflecting more stable semantic inheritance. These trajectories exemplify how different T2I architectures encode, preserve, or distort textual intent, offering a geometry-grounded framework for auditing \textbf{generative alignment}, \textbf{hallucination risk}, and \textbf{referential fidelity} in text-to-image synthesis.
    }
    \label{fig:ndna_t2i_trajectory}
\end{figure}



\clearpage

\section*{Glimpse from Ongoing Work on nDNAs of Embodied Foundation Models (EFMs)}

\vspace{-2mm}
\noindent
This section offers an \emph{early glimpse} into our ongoing efforts to characterize the \textbf{neural genomic structures} (\textbf{nDNAs}) of \textbf{EFMs}. Unlike vision-language models or pure language models, EFMs fuse perception, reasoning, and action across layers -- generating complex epistemic pathways that standard scalar metrics fail to capture. Our proposed \textbf{nDNA-EFM} framework provides a principled, geometry-grounded signature of these pathways, exposing how semantics, policies, and alignment dynamics evolve as latent representations traverse the perception-to-action stack.  

Each nDNA trajectory integrates multiple diagnostic components: \textbf{Spectral Curvature} ($\kappa^{(l)}$) measures latent manifold entanglement, highlighting ambiguous grounding or policy bifurcations; \textbf{Thermodynamic Length} ($\mathcal{L}^{(l)}$) quantifies accumulated semantic displacement across layers, signaling epistemic effort or overcompression; \textbf{Latent Torsion} ($\tau^{(l)}$) captures non-planarity of the latent trajectory, flagging latent drift between perception and intent; \textbf{Policy Vector Field Magnitude} reflects the directional sensitivity of the policy to latent shifts; \textbf{Policy Torsion Index} detects inconsistency or directional twist in action selection; and \textbf{Lie Derivative Equivariance Error} measures deviation from symmetry-preserving policy dynamics.  

Preliminary results (see Figure~\ref{fig:ndna_efm_trajectory}) demonstrate how \textbf{nDNA-EFM} moves beyond surface-level output metrics, uncovering the \emph{hidden geometry of knowledge, intent, and policy evolution} in embodied AI. This framework offers a fundamentally new lens for \textbf{alignment auditing}, \textbf{policy interpretability}, and \textbf{epistemic traceability}, revealing not just what decisions EFMs make, but how those decisions emerge through their latent dynamics.



\begin{figure}[htbp]
    \centering
    \includegraphics[width=0.58\textwidth]{Figures/ndna_efm_trajectories.png}
    \caption{
    \textbf{nDNA Trajectories for Embodied Foundation Models (EFMs): Geometric Evolution in Latent Policy Space.}
    This 3D plot visualizes the \textbf{neural DNA} (\textbf{nDNA}) trajectories of five state-of-the-art EFMs -- \textbf{RT-2}~\citep{brohan2023rt2}, \textbf{PaLM-E}~\citep{driess2023palme}, \textbf{SayCan}~\citep{ahn2022saycan}, \textbf{VIMA}~\citep{jiang2022vima}, and \textbf{GR00T}~\citep{singh2023gr00t} -- across their final transformer layers. The latent manifold is spanned by three principal axes: \emph{Spectral Curvature} ($\kappa^{(l)}$), reflecting local semantic entanglement; \emph{Thermodynamic Length} ($\mathcal{L}^{(l)}$), quantifying epistemic “semantic effort” across layers; and \emph{Latent Torsion} ($\tau^{(l)}$), capturing directional twist between evolving policy and perception embeddings.  
    \textbf{RT-2} and \textbf{PaLM-E} trace planar arcs with low torsion and stable curvature, suggesting coherent instruction-following and well-aligned visual grounding. \textbf{SayCan} shows extended thermodynamic length and moderate curvature, consistent with its hierarchical plan decomposition and constrained decision logic. \textbf{VIMA} exhibits episodic torsion spikes, reflecting latent twists from aggressive zero-shot generalization across synthetic scenes. \textbf{GR00T} displays pronounced curvature and torsion fluctuations, indicative of complex perception-action fusion in high-fidelity 3D simulation, and epistemic strain during policy blending.  
    These trajectories demonstrate how \textbf{nDNA-EFM} reveals \emph{how models arrive at action}, not merely \emph{what they do}. Unlike scalar metrics, this geometric lens exposes the \textbf{path-dependence of knowledge, intent}, and \textbf{alignment dynamics} in embodied AI.
    }
    \label{fig:ndna_efm_trajectory}
\end{figure}



\clearpage

\section*{Glimpse from Ongoing Work on nDNAs of Diffusion LMs}

\vspace{-2mm}
\noindent
This section offers an \emph{early glimpse} into our ongoing study of the \textbf{neural genomic structures} (\textbf{nDNAs}) of \textbf{diffusion-based language models} (\textbf{Diffusion LMs}). Unlike autoregressive large language models (LLMs) that extend text token-by-token~\citep{brown2020gpt3, touvron2023llama, chowdhery2022palm}, Diffusion LMs such as \textbf{Genie}~\citep{genie2024openlm} and \textbf{CodeGenie}~\citep{codegenie2024openlm} synthesize language by iteratively denoising latent embeddings -- a paradigm inspired by successes in vision and multimodal diffusion~\citep{ho2020denoising, saharia2022photorealistic, nichol2021glide}. This fundamental difference raises questions about their latent semantic scaffolds, representational complexity, and capacity for alignment and compositional generalization~\citep{li2022diffusiondb, yang2023diffusionlm, ge2023diffusionllm}.  

Our \textbf{nDNA-Diffusion} framework maps these models into a latent geometry characterized by \emph{Spectral Curvature} ($\kappa_t$), measuring manifold complexity~\citep{belkin2003laplacian, coifman2006diffusion}; \emph{Thermodynamic Length} ($\mathcal{L}_t$), capturing accumulated semantic displacement~\citep{crooks2007measuring, sivak2012thermodynamic}; and \emph{Diffusion Time} ($t$), indexing generative progression across denoising steps~\citep{ho2020denoising, nichol2021glide}. Figure~\ref{fig:ndna_diffusion_trajectories} illustrates how diffusion language models like \textbf{Genie} and \textbf{CodeGenie} differ in their latent epistemic pathways, despite producing fluent outputs. The nDNA trajectories reveal contrasts in semantic curvature, thermodynamic length, and belief evolution that are otherwise invisible in surface-level performance.


Preliminary comparisons (see Figure~\ref{fig:ndna_diffusion_vs_llama}) show that Diffusion LMs trace \emph{flatter}, more compressed latent trajectories relative to autoregressive models like \textbf{LLaMA}~\citep{touvron2023llama}, whose nDNA reveals richer curvature and layered semantic buildup -- patterns associated with deeper compositional reasoning~\citep{wei2022chain, nye2021improving} and robust alignment~\citep{ouyang2022training, ganguli2023reducing}. These early results align with prior suggestions that diffusion LMs, while promising~\citep{ge2023diffusionllm}, are still in their \emph{epistemic infancy}, lacking the semantic scaffold complexity seen in mature LLMs.  

Our ongoing work seeks to position \textbf{nDNA geometry} as a fundamentally new lens for auditing the \emph{alignment dynamics} of generative models~\citep{perez2022discovering}, probing their \emph{epistemic maturity} and capacity for robust reasoning~\citep{bommasani2021opportunities}, and benchmarking the \emph{developmental trajectories} that distinguish diffusion-based language models from their autoregressive counterparts. By moving beyond output-level metrics, this approach aims to reveal the latent structural pathways through which models encode, transform, and act on knowledge.

\begin{figure*}[htbp]
    \centering
    \includegraphics[width=0.95\linewidth]{diffusion/diffusion_2_models.png}
    \caption{
    \textbf{nDNA Trajectories for OpenLM-Based Diffusion Language Models.}
    This figure illustrates the \textbf{neural DNA (nDNA)} trajectories of two diffusion-based language models -- \textbf{Genie (OpenLM)}~\citep{genie2024openlm} and \textbf{CodeGenie (OpenLM)}~\citep{codegenie2024openlm} -- within a latent geometry defined by three core components: \textbf{Thermodynamic Length} ($\mathcal{L}_t$), capturing accumulated semantic displacement; \textbf{Spectral Curvature} ($\kappa_t$), reflecting latent manifold complexity; and \textbf{Diffusion Time} ($t$), denoting progression along the generative process.  
    Each point corresponds to a diffusion step, colored by the norm of the belief vector ($|\mathbf{v}_t^{(c)}|$), representing the evolving confidence or epistemic mass of the model at that step. \textbf{Genie} exhibits a smoother, more uniform latent arc with moderate curvature, indicating stable semantic evolution during generation. In contrast, \textbf{CodeGenie} shows regions of sharper curvature and concentrated belief mass at early diffusion times, reflecting more abrupt semantic reconfiguration during code synthesis.  
    These trajectories highlight how \textbf{nDNA-Diffusion} geometry reveals differences in epistemic flow, semantic complexity, and alignment dynamics not visible from final outputs alone. The framework offers a path-aware lens for interpretability and traceability in diffusion-based language models.
    }
    \label{fig:ndna_diffusion_trajectories}
\end{figure*}



\begin{figure}[htbp]
    \centering
    \includegraphics[width=0.7\linewidth]{diffusion/diffusion_lm_vs_llama.png}
    \caption{
    \textbf{nDNA Comparison: LLaMA vs Diffusion Language Models (Genie, CodeGenie).}
    This 3D plot compares the \textbf{neural DNA (nDNA)} trajectories of an autoregressive large language model (\textbf{LLaMA}) and two diffusion-based language models -- \textbf{Genie} and \textbf{CodeGenie} -- at a shared latent geometric scale. The axes represent \emph{Spectral Curvature} ($\kappa_t$), \emph{Thermodynamic Length} ($\mathcal{L}_t$), and \emph{Layer/Step Index} ($t$).  
    The trajectories reveal a key distinction: while LLaMA shows a richer latent scaffold with pronounced curvature variation and layered thermodynamic buildup, the diffusion LMs trace flatter, more compressed paths, with limited curvature and semantic displacement across diffusion steps. This suggests that \textbf{diffusion language models may still be at an early stage of epistemic maturity}, with latent scaffolds that are less complex than those of established autoregressive LLMs. 
    }
    \label{fig:ndna_diffusion_vs_llama}
\end{figure}
%\clearpage



\section*{\emph{n}DNA v2 -- Tensor-Rich Belief Torsion}
\label{sec:ndna_tensor}

While our current framework encapsulates belief as a vector--valued semantic steering force, it is possible to \textbf{extend this formulation} to a richer \emph{tensor-based representation} of belief, shown in Table~\ref{tab:ndna_tensor_example}.

This higher-order representation \(\mathbf{T}_\ell^{(c)} \in \mathbb{R}^{d \times d \times \cdots}\) captures \emph{complex multi-dimensional interactions} within the latent activation space, enabling the modeling of \textbf{nonlinear dependencies}, \emph{higher-order correlations}, and \textbf{contextual entanglements} that are fundamentally inaccessible to vector embeddings alone. By encoding semantic \emph{torsion}, \emph{curvature}, and \emph{multi-modal interactions} across layers, tensor-based belief representation offers a powerful framework to characterize subtle latent geometry and richer internal dynamics of language models.

Such a tensor-rich approach is particularly beneficial in cases involving \textbf{multimodal fusion}, \emph{hierarchical reasoning}, and \textbf{complex compositionality}, where simple vectorial steering fails to capture \emph{intertwined semantic constraints} and \emph{latent contextual shifts}. For example, in tasks demanding \emph{disentanglement of overlapping concepts} or \textbf{fine-grained alignment} of multi-turn dialogue states, tensors can explicitly encode relational structure and layered semantic dependencies.

However, this expressivity comes with a significant \textbf{computational overhead}, as tensor operations increase dimensionality exponentially and necessitate \emph{efficient approximation schemes} or \emph{structured sparsity} to remain scalable. Balancing the tradeoff between \textbf{modeling fidelity} and \textbf{computational tractability} is a central challenge we are actively addressing in ongoing research, with promising directions including \emph{low-rank tensor decompositions} and \emph{hybrid vector-tensor architectures}.

In summary, while belief-as-tensor unlocks \emph{richer representational nuance} and \textbf{deeper semantic insight}, its practical deployment requires careful design to harness its power without prohibitive costs - we are currently working.






\begin{table}[H]
\centering
\scriptsize
\caption{
\textbf{Illustrative nDNA Tensor Example:} Presents a combined view of \textbf{spectral curvature} \(\boldsymbol{\kappa_\ell}\), \textbf{thermodynamic length} \(\boldsymbol{\mathcal{L}_\ell}\), \textbf{Frobenius norm of tensor slices} \(\|\mathbf{T}_\ell^{(c)}\|_F\), alongside \textbf{visual heatmaps} of the corresponding tensor, highlighting the intricate latent geometry at different layers.
}
\label{tab:ndna_tensor_example}
\vspace{-0.5em}
\begin{tabular}{|c|c|c|c|c|}
\hline
\textbf{Layer} 
& \(\boldsymbol{\kappa_\ell}\) 
& \(\boldsymbol{\mathcal{L}_\ell}\) 
& \(\|\mathbf{T}_\ell^{(c)}\|_F\) 
& \textbf{Tensor Heatmap} \\
\hline
20 & \cellcolor{green!20}0.0412 & \cellcolor{yellow!30}0.9123 & \cellcolor{orange!20}1.4521 
& \raisebox{-0.5\height}{\includegraphics[width=2.5cm]{tensor_heatmaps/layer_20_tensor.png}} \\
21 & \cellcolor{green!30}0.0458 & \cellcolor{green!15}0.8123 & \cellcolor{red!30}1.7523 
& \raisebox{-0.5\height}{\includegraphics[width=2.5cm]{tensor_heatmaps/layer_21_tensor.png}} \\
\hline
\multicolumn{5}{|c|}{\(\vdots\)} \\
\hline
27 & \cellcolor{yellow!20}0.0512 & \cellcolor{red!30}1.0520 & \cellcolor{orange!30}1.6222 
& \raisebox{-0.5\height}{\includegraphics[width=2.5cm]{tensor_heatmaps/layer_27_tensor.png}} \\
\hline
\multicolumn{5}{|c|}{\(\vdots\)} \\
\hline
29 & \cellcolor{red!30}0.0672 & \cellcolor{orange!30}1.0123 & \cellcolor{yellow!20}1.4322 
& \raisebox{-0.5\height}{\includegraphics[width=2.5cm]{tensor_heatmaps/layer_29_tensor.png}} \\
30 & \cellcolor{orange!20}0.0555 & \cellcolor{green!15}0.8221 & \cellcolor{red!20}1.6720 
& \raisebox{-0.5\height}{\includegraphics[width=2.5cm]{tensor_heatmaps/layer_30_tensor.png}} \\
\hline
\end{tabular}
\end{table}
